% Generated draft: requires mathematical and source review.
\documentclass[11pt]{article}
\usepackage[margin=25mm]{geometry}
\usepackage{fontspec}
\setmainfont{Latin Modern Roman}
\usepackage{amsmath,amssymb,mathtools,unicode-math}
\setmathfont{Latin Modern Math}
\usepackage{enumitem,needspace,xurl,hyperref}
\setlength{\parindent}{0pt}
\setlength{\parskip}{5pt}
\setlength{\emergencystretch}{3em}
\newcommand{\N}{\mathbb{N}}
\newcommand{\Z}{\mathbb{Z}}
\newcommand{\Q}{\mathbb{Q}}
\newcommand{\R}{\mathbb{R}}
\newcommand{\C}{\mathbb{C}}
\newcommand{\F}{\mathbb{F}}
\newcommand{\A}{\mathbb{A}}
\newcommand{\PP}{\mathbb{P}}
\newcommand{\E}{\mathbb{E}}
\newcommand{\eps}{\varepsilon}
\newcommand{\dd}{\,\mathrm{d}}
\begin{document}
\section*{UnsolvedMath ID: 11000293}
\textbf{Source problem code:} catalogue entry
\subsection{Definitions and mathematical statement}
Let \(F_n\) be the rank-\(n\) free group. Extending automorphisms by fixing a new generator gives \(s_n:\operatorname{Aut}(F_n)\to\operatorname{Aut}(F_{n+1})\). The problem asks for the exact rational homological stability threshold as a function of degree \(i\): the least rank after which the corresponding homology is stable, both for automorphism and outer automorphism groups.

For automorphisms this can be expressed as the least \(N\) such that every \((s_n)_*\) on \(H_i(-;\Q)\) is an isomorphism for \(n\geq N\). The supplied summary reports Galatius's stable rational homology as zero in positive degrees. In that interpretation, a map-independent version for either group family \(G_n\) is
\[
 N_G(i)=\min\{N:H_i(G_n;\Q)=0\text{ for every }n\geq N\},
 \qquad i\geq1.
\]
Degree zero is \(\Q\), not zero. The threshold is not the first isolated vanishing rank. For \(\operatorname{Out}(F_n)\), the source's comparison/stabilization model is missing: basis extension on automorphisms does not simply descend to a canonical homomorphism on outer automorphisms, so such a map is not invented.

The source reports only two nontrivial unstable outer-group classes then known [49, 26]. It describes Morita's infinite series of cycles [69], using Kontsevich's connection with cohomology of an infinite-dimensional Lie algebra. The first generates
\(H_4(\operatorname{Out}(F_4);\Q)\cong\Q\), and Conant--Vogtmann [26] prove nontriviality of the second in \(H_8(\operatorname{Out}(F_6);\Q)\). Morita and Conant--Vogtmann also construct cycles indexed by odd-valent graphs. Cycle construction alone is not asserted to make every member a nonzero class. The exact Lie-algebra/weight correspondence is not supplied.
\subsection{Short English statement}
Determine the optimal rational stability or positive-degree vanishing rank for each homological degree of free-group automorphism and outer automorphism groups, retaining the reported unstable Morita classes.
\subsection{Sources}
\begin{itemize}
\item \emph{Problems on Mapping Class Groups and Related Topics}, edited by Farb, Question 5.1 and continuation, supplied PDF page 336. \url{https://www.math.uchicago.edu/~farb/papers/mcgbook.pdf}. Recorded access 2026-07-29; extraction pdf-font-structure.
\item [49, 26] for the then-known unstable classes; Morita [69]; Conant and Vogtmann [26]; Kontsevich's Lie-algebra correspondence. Complete bibliography and precise comparison statements are absent. The summary credits Galatius with positive-degree stable rational vanishing but supplies no full reference.
\item Local UnsolvedMath registry, \texttt{data/unsolvedmath/problems.json}, permanent ID \(11000293\), exact source code \texttt{catalogue entry}; supplied OPEN-TRIAGE summary distinguishes known linear ranges from unknown exact thresholds.
\end{itemize}
\end{document}
